% Source: https://artofproblemsolving.com/community/c3h1762370_lowest__number
% Stable identity: aops-topic-1762370-q01
\textbf{Number of digits.} Let the number have $m$ nonzero digits (and one $0$). They include $1,2,\dots,8$ (sum $36$) and $m-8$ further digits, each at most $9$, so $36+9(m-8)\ge2018$, i.e.\ $9(m-8)\ge1982$, giving $m\ge229$. A number with fewer digits is smaller, so we look for the smallest number with $229$ nonzero digits and one zero, i.e.\ $230$ digits, if one exists.

\textbf{Which digits.} With $229$ nonzero digits, the total shortfall from all nines is $9\cdot229-2018=43$. The obligatory digits $1,\dots,8$ account for a shortfall of $8+7+\cdots+1=36$, so the $221$ other nonzero digits have total shortfall $7$ (a digit $d$ has shortfall $9-d$).

\textbf{Smallest arrangement.} The first digit cannot be $0$, so it is at least $1$; we can make it $1$. The second digit can be $0$. A second $1$ would have shortfall $8>7$, so there is exactly one $1$ and the third digit is at least $2$; it can be $2$. The fourth digit can be $2$ only if one of the other digits is a $2$, whose shortfall $7$ uses up the whole allowance; this forces the digits to be $0,1,2,2,3,4,5,6,7,8$ and $220$ nines. Any other choice makes the fourth digit at least $3$ and the number larger. With this multiset the smallest arrangement puts $1$ first, then $0$, then the rest in increasing order:
\[
\boxed{10223456789\cdots9}\quad(\text{the digits }1,0,2,2,3,4,5,6,7,8\text{ followed by }220\text{ nines}).
\]
Check: $1+0+2+2+3+4+5+6+7+8=38$ and $38+220\cdot9=2018$; it has $230$ digits.
